\subsection{层级（Echelons）}

一个层级构造方案是自然整数有序对的一个序列 \(c_1, c_2, \ldots, c_m\) （*） \(c_i = (a_i, b_i)\) ，满足以下条件：

\begin{enumerate}
\def\labelenumi{(\alph{enumi})}
\item
  如果 \(b_{i} = 0\) ，那么 \(1 \leqslant a_{i} \leqslant i - 1\) .
\item
  如果 \(a_i \neq 0\) 并且 \(b_i \neq 0\) ，那么 \(1 \leqslant a_i \leqslant i - 1\) 并且 \(1 \leqslant b_i \leqslant i - 1\) .
\end{enumerate}

这些条件蕴含 \(c_{1}=(0,b_{1})\) ，其中 \(b_{1}>0\) 。如果 n 是整数 \(b_{i}\) （它们出现在对 \((0,b_{i})\) 中）中最大的一个，那么 \(c_{1},c_{2},\ldots,c_{m}\) 则称其为 n 项上的层级构造方案。

给定一个 n 项上的层级构造方案 \(\mathbf{S} = (c_{1}, c_{2}, \ldots, c_{m})\) ，并给定理论 C 中的 n 个项 \(E_{1}, E_{2}, \ldots, E_{n}\) ，其中 C 是比集合论更强的理论，则方案 S 在 \(E_{1}, \ldots, E_{n}\) 上的一个层级构造被定义为理论 C 中的一个 m 项序列 \(A_{1}, A_{2}, \ldots, A_{m}\) ，由以下条件逐步定义：

\begin{enumerate}
\def\labelenumi{(\alph{enumi})}
\item
  如果 \(c_{i} = (0, b_{i})\) ，那么 \(A_{i}\) 是项 \(E_{b_{i}}\) .
\item
  如果 \(c_{i} = (a_{i}, 0)\) ，那么 \(A_{i}\) 是项 \(\mathfrak{P}(A_{a_i})\) .
\item
  如果 \(c_{i} = (a_{i}, b_{i})\) ，其中 \(a_{i} \neq 0\) 并且 \(b_{i} \neq 0\) ，那么 \(A_{i}\) 是项 \(A_{a_{i}} \times A_{b_{i}}\) .
\end{enumerate}

最后一项 \(A_{m}\) ，即方案 S 在 \(E_{1}, \ldots, E_{n}\) 上的阶梯构造的末项，称为方案 S 在基集合 \(E_{1}, \ldots, E_{n}\) 上的阶梯；在下文的一般论证中，它记作 \(S(E_{1}, \ldots, E_{n})\) .

示例。给定两个集合E，F，集合 \(\mathfrak{P}(\mathfrak{P}(E)) \times \mathfrak{P}(F)\) 是E，F上的一个阶梯，具有概形

\[
(0, 1), (0, 2), (1, 0), (3, 0), (2, 0), (4, 5).
\]

它也是E，F上的梯队，具有方案

\[
(0, 2), (0, 1), (1, 0), (2, 0), (4, 0), (5, 3).
\]

因此，不同的方案可能会在相同的条件下产生相同的梯队。

